% archon:covers MR2223407ConstructionHilbertQuot/Variants.lean
% Variants and applications of the Quot existence theorem (Nitsure §6):
% the quasi-projective Quot scheme, the Hilbert scheme, the scheme of
% morphisms, and the quotient by a flat projective equivalence relation.
% Source: Nitsure, "Construction of Hilbert and Quot Schemes" (MR2223407,
% arXiv:math/0504590), file references/MR2223407-construction-hilbert-quot.tex.

\chapter{Variants and Applications}
\label{chap:variants}

This chapter extends the main Quot existence theorem
(\cref{thm:altman-kleiman-quot}) from the projective to the quasi-projective
setting, and derives the three headline applications: the Hilbert scheme, the
scheme of morphisms, and the quotient of a scheme by a flat projective
equivalence relation. Every result here is a consequence of the goal theorem
and of the material developed in the earlier chapters; none re-proves the Quot
construction itself.

\section{Quot scheme in the quasi-projective case}

The passage from the projective to the quasi-projective case rests on a single
openness fact: restricting a Quot functor to an open subscheme of the ambient
space cuts out an open subfunctor. This is in turn a consequence of an
elementary statement about the locus over which the support of a coherent sheaf
avoids a fixed closed subscheme.

\begin{lemma}\leanok
[Quot of an open is an open subfunctor]
  \label{lem:quot-open-subfunctor}
  \lean{MR2223407ConstructionHilbertQuot.quot_open_subfunctor}
  \uses{def:quot-functor}
  % SOURCE: Nitsure, §6, lines 2615--2634 (read from references/MR2223407-construction-hilbert-quot.tex)
  % SOURCE QUOTE: "\textbf{Exercise }
  % Let $\pi : Z\to S$ be a proper morphism of noetherian schemes.
  % Let $Y\subset Z$ be a closed subscheme, and let $\F$ be a coherent sheaf
  % on $Z$. Then there exists an open subscheme $S' \subset S$
  % with the universal property that a morphism $T\to S$ factors through
  % $S'$ if and only if the support of the pull-back
  % $\F_T$ on $Z_T = Z\times_ST$ is disjoint from
  % $Y_T = Y \times_ST$.
  % \textbf{Exercise } As a consequence of the above, show the following:
  % If $\pi : Z\to S$ is a proper morphism with $S$ noetherian, if
  % $X\subset Z$ is an open subscheme, and if
  % $E$ is a coherent sheaf on $Z$, then
  % $\Quot_{E|_X/X/S}$ is an open subfunctor of $\Quot_{E/Z/S}$."
  \textit{Source: Nitsure §6 (two exercises).}
  Let $\pi\colon Z\to S$ be a proper morphism of noetherian schemes, $Y\subset Z$
  a closed subscheme and $\Ff$ a coherent sheaf on $Z$. Then there is an open
  subscheme $S'\subset S$ with the universal property that a morphism $T\to S$
  factors through $S'$ if and only if the support of the pull-back $\Ff_T$ on
  $Z_T = Z\times_S T$ is disjoint from $Y_T = Y\times_S T$. As a consequence, if
  $X\subset Z$ is an open subscheme and $E$ is a coherent sheaf on $Z$, then
  $\Quotf_{\restr{E}{X}/X/S}$ is an open subfunctor of $\Quotf_{E/Z/S}$.
\end{lemma}

\begin{proof}
  \uses{def:quot-functor}
  % This statement is posed as two exercises in Nitsure; the proof below is the
  % standard argument (no verbatim source proof is available to quote).
  For the first assertion, the support $\mathrm{Supp}\,\Ff$ is a closed subset of
  $Z$, hence so is $W = \mathrm{Supp}\,\Ff\cap Y$. Since $\pi$ is proper, the image
  $\pi(W)$ is closed in $S$; set $S' = S\setminus \pi(W)$ with its open subscheme
  structure. Formation of the support of a coherent sheaf commutes with the flat
  (indeed arbitrary, on the level of underlying sets) base change $Z_T\to Z$, so
  $\mathrm{Supp}\,\Ff_T = \mathrm{Supp}(\Ff)\times_S T$; likewise
  $Y_T = Y\times_S T$. Hence $\mathrm{Supp}\,\Ff_T$ is disjoint from $Y_T$ exactly
  when the image of $T\to S$ avoids $\pi(W)$, i.e.\ when $T\to S$ factors through
  $S'$. This is the required universal property.

  For the consequence, write $Y = Z\setminus X$ with its reduced (or any) closed
  structure. A $T$-valued point of $\Quotf_{E/Z/S}$ is a coherent quotient
  $q\colon E_T\twoheadrightarrow \Gg$ on $Z_T$ with $\Gg$ flat over $T$ and of
  support proper over $T$; it lies in the subfunctor $\Quotf_{\restr{E}{X}/X/S}$
  precisely when $\mathrm{Supp}\,\Gg$ is disjoint from $Y_T$ (so that the quotient
  is supported in $X_T$ and is determined by its restriction to $X_T$). Applying
  the first assertion to the proper morphism $Z_T\to T$, the closed subscheme
  $Y_T$ and the coherent sheaf $\Gg$ shows that the locus of base changes
  satisfying this condition is represented by an open subscheme of $T$. Thus the
  inclusion of functors is representable by open immersions, i.e.\
  $\Quotf_{\restr{E}{X}/X/S}$ is an open subfunctor of $\Quotf_{E/Z/S}$.
\end{proof}

\begin{theorem}\leanok
[Strongly quasi-projective Quot, Altman--Kleiman]
  \label{thm:strongly-quasiprojective-quot}
  \lean{MR2223407ConstructionHilbertQuot.stronglyQuasiprojectiveQuot}
  \uses{lem:quot-open-subfunctor, thm:altman-kleiman-quot, def:strongly-projective, def:represents-functor}
  % SOURCE: Nitsure, §6, lines 2645--2657 (read from references/MR2223407-construction-hilbert-quot.tex)
  % SOURCE QUOTE: "\begin{theorem}\label{Strongly quasi-projective quot}
  % {\rm (Altman and Kleiman)}
  % Let $S$ be a noetherian scheme, $X$ a locally closed subscheme of
  % $\PP(V)$ for some vector bundle $V$ on $S$, $L = {\mathcal O}_{\PP(V)}(1)|_X$,
  % $E$ a coherent quotient sheaf of $\pi^*(W)(\nu)|_X$ where
  % $W$ is a vector bundle on $S$ and $\nu$ is an integer, and
  % $\Phi \in \Q[\lambda]$.
  % Then the functor $\Quot_{E/X/S}^{\Phi,L}$ is representable
  % by a scheme $\quot_{E/X/S}^{\Phi,L}$ which can be embedded over $S$ as a
  % locally closed subscheme of $\PP(F)$ for some vector bundle $F$ on $S$.
  % Moreover, the vector bundle $F$ can be chosen to be an exterior power of
  % the tensor product of $W$ with a symmetric power of $V$.
  % \end{theorem}"
  \textit{Source: Nitsure §6, Theorem (Strongly quasi-projective quot).}
  Let $S$ be a noetherian scheme, $X$ a locally closed subscheme of $\PP(V)$ for
  a vector bundle $V$ on $S$, $L = \restr{\Oo_{\PP(V)}(1)}{X}$, $E$ a coherent
  quotient sheaf of $\restr{\pi^*(W)(\nu)}{X}$ where $W$ is a vector bundle on $S$
  and $\nu\in\ZZ$, and $\Phi\in\QQ[\lambda]$. Then $\Quotf_{E/X/S}^{\Phi,L}$ is
  representable by a scheme $\quotsch_{E/X/S}^{\Phi,L}$ which embeds over $S$ as a
  locally closed subscheme of $\PP(F)$ for some vector bundle $F$ on $S$; $F$ may
  be taken to be an exterior power of the tensor product of $W$ with a symmetric
  power of $V$.
\end{theorem}

% SOURCE QUOTE PROOF: "\proof Let $\ov{X} \subset \PP(V)$ be the schematic closure of
% $X\subset \PP(V)$, and let $\ov{E}$ be the coherent sheaf on $\ov{X}$
% defined as the image of the composite homomorphism
% $$\pi^*(W)(\nu)|_{\ov{X}} \to j_*(\pi^*(W)(\nu)|_X) \to j_*E$$
% Then we get a quotient $\pi^*(W)(\nu)|_{\ov{X}} \to \ov{E}$
% which restricts on $X\subset \ov{X}$ to the given quotient
% $\pi^*(W)(\nu)|_X\to E$. Therefore by the above exercise,
% $\Quot_{E/X/S}$ is an open subfunctor of
% $\Quot_{\ov{E}/\ov{X}/S}$. Now the result follows from the
% Theorem \ref{Altman-Kleiman Quot}.\qed"
\begin{proof}
  \uses{lem:quot-open-subfunctor, thm:altman-kleiman-quot}
  Let $\overline X\subset \PP(V)$ be the schematic closure of $X$, an open
  subscheme $j\colon X\hookrightarrow \overline X$, and let $\overline E$ be the
  image of the composite
  $\restr{\pi^*(W)(\nu)}{\overline X}\to j_*(\restr{\pi^*(W)(\nu)}{X})\to j_*E$.
  This gives a quotient $\restr{\pi^*(W)(\nu)}{\overline X}\to \overline E$ whose
  restriction to $X$ is the given quotient $\restr{\pi^*(W)(\nu)}{X}\to E$. By
  \cref{lem:quot-open-subfunctor}, $\Quotf_{E/X/S}$ is an open subfunctor of
  $\Quotf_{\overline E/\overline X/S}$. Since $\overline X$ is a closed subscheme
  of $\PP(V)$, hence strongly projective over $S$, \cref{thm:altman-kleiman-quot}
  represents $\Quotf_{\overline E/\overline X/S}^{\Phi,L}$ by a closed subscheme
  of a suitable $\PP(F)$; the open subfunctor $\Quotf_{E/X/S}^{\Phi,L}$ is then
  represented by the corresponding open part, a locally closed subscheme of
  $\PP(F)$.
\end{proof}

To extend Grothendieck's construction to the quasi-projective case one needs the
following prolongation lemma, of independent interest.

\begin{lemma}\leanok
[Coherent prolongation]
  \label{lem:coherent-prolongation}
  \lean{MR2223407ConstructionHilbertQuot.coherent_prolongation}
  % SOURCE: Nitsure, §6, lines 2679--2682 (read from references/MR2223407-construction-hilbert-quot.tex)
  % SOURCE QUOTE: "\begin{lemma}\label{coherent prolongation}
  % Any coherent sheaf on an open subscheme of a noetherian scheme $S$ can be
  % prolonged to a coherent sheaf on all of $S$.
  % \end{lemma}"
  \textit{Source: Nitsure §6, Lemma (coherent prolongation).}
  Any coherent sheaf on an open subscheme of a noetherian scheme $S$ can be
  prolonged to a coherent sheaf on all of $S$.
\end{lemma}

% SOURCE QUOTE PROOF: "\proof First consider the case where $S = \spec A$ is affine, and
% let $j: U\hra S$ denote the inclusion.
% The quasi-coherent sheaf
% $j_*(\F)$ corresponds to the $A$-module $M = H^0(S,j_*(\F))$,
% in the sense that $j_*(\F) = M^{\sim}$.
% Given any $u\in U$, there exist finitely many
% elements $e_1,\ldots,e_n \in M$
% which generate the fiber $\F_u$ regarded as a vector space
% over the residue field $\kappa(u)$. By Nakayama these elements
% will generate the stalks of $\F$ in an open neighbourhood of $u$ in $U$.
% Therefore by the noetherian hypothesis, there exist finitely many
% elements $e_1,\ldots,e_r \in M$ which generate the stalk of $\F$
% at each point of $U$. If $N\subset M$ is the submodule generated by these
% elements, then $\G = N^{\sim}$ is a coherent prolongation of
% $\F$ to $S = \spec A$, proving the result in the affine case.
% In the general case, by the
% noetherian condition there exists a maximal coherent prolongation
% $(U',\F')$ of $\F$. Then unless $U' = S$, we can obtain
% a further prolongation of $\F'$ by using the affine case.
% For, if $u\in S- U'$,
% we can take an affine open subscheme $V$ containing $u$, and a
% coherent prolongation $\G'$ of $\F'|_{U'\bigcap V}$ to all of $V$,
% and then glue together $\G'$ and $\F'$ along $U'\bigcap V$ to
% further prolong $\F'$ to $U'\bigcup V$,
% contradicting the maximality of $(U',\F')$.
% \qed"
\begin{proof}
  First let $S = \spec A$ be affine with inclusion $j\colon U\hookrightarrow S$.
  The quasi-coherent sheaf $j_*\Ff$ corresponds to the $A$-module
  $M = H^0(S,j_*\Ff)$, i.e.\ $j_*\Ff = \widetilde M$. For each $u\in U$ pick
  finitely many $e_1,\dots,e_n\in M$ generating the fibre $\Ff_u$ over
  $\kappa(u)$; by Nakayama these generate the stalk of $\Ff$ on an open
  neighbourhood of $u$ in $U$, and by noetherianity finitely many
  $e_1,\dots,e_r\in M$ generate the stalk of $\Ff$ at every point of $U$. If
  $N\subset M$ is the submodule they generate, then $\Gg = \widetilde N$ is a
  coherent prolongation of $\Ff$ to $S$. In general, by noetherianity there is a
  maximal coherent prolongation $(U',\Ff')$ of $\Ff$. If $U'\ne S$, choose
  $u\in S\setminus U'$ and an affine open $V\ni u$; the affine case prolongs
  $\restr{\Ff'}{U'\cap V}$ to $V$, and gluing along $U'\cap V$ extends $\Ff'$ to
  $U'\cup V$, contradicting maximality. Hence $U' = S$.
\end{proof}

\begin{theorem}\leanok
[Quasi-projective Quot, Grothendieck]
  \label{thm:quasiprojective-quot}
  \lean{MR2223407ConstructionHilbertQuot.quasiprojectiveQuot}
  \uses{lem:coherent-prolongation, lem:quot-open-subfunctor, thm:grothendieck-quot, def:represents-functor}
  % SOURCE: Nitsure, §6, lines 2713--2721 (read from references/MR2223407-construction-hilbert-quot.tex)
  % SOURCE QUOTE: "\begin{theorem}\label{Quasi-projective quot}{\rm (Grothendieck)}
  % Let $S$ be a noetherian scheme, $X$ a quasi-projective
  % scheme over $S$, $L$ a line bundle on $X$ which is
  % relatively very ample over $S$, $E$ a quotient sheaf on $X$,
  % and $\Phi \in \Q[\lambda]$.
  % Then the functor $\Quot_{E/X/S}^{\Phi,L}$ is representable
  % by a scheme $\quot_{E/X/S}^{\Phi,L}$ which is
  % quasi-projective over $S$.
  % \end{theorem}"
  \textit{Source: Nitsure §6, Theorem (Quasi-projective quot).}
  Let $S$ be a noetherian scheme, $X$ a quasi-projective scheme over $S$, $L$ a
  relatively very ample line bundle on $X$, $E$ a quotient sheaf on $X$ and
  $\Phi\in\QQ[\lambda]$. Then $\Quotf_{E/X/S}^{\Phi,L}$ is representable by a
  scheme $\quotsch_{E/X/S}^{\Phi,L}$ which is quasi-projective over $S$.
\end{theorem}

% SOURCE QUOTE PROOF: "\proof By definition of quasi-projectivity of $X\to S$, note that
% $X$ can be embedded over $S$ as a locally closed subscheme of
% $\PP(V)$ for some coherent sheaf $V$ on $S$,
% such that $L$ is isomorphic to ${\mathcal O}_{\PP(V)}(1)|_X$.
% Let $\ov{X} \subset \PP(V)$ be the schematic closure of
% $X$ in $\PP(V)$. This is a projective scheme over $S$, and
% $X$ is embedded as an open subscheme in it.
% By Lemma \ref{coherent prolongation} the
% coherent sheaf $E$ has a coherent prolongation $\ov{E}$ to $\ov{X}$.
% For any such prolongation $\ov{E}$, the functor
% $\Quot_{E/X/S}$ is an open subfunctor of $\Quot_{\ov{E}/\ov{X}/S}$.
% Therefore the desired result now follows from Theorem \ref{Grothendieck Quot}.
% \qed"
\begin{proof}
  \uses{lem:coherent-prolongation, lem:quot-open-subfunctor, thm:grothendieck-quot, def:represents-functor}
  By quasi-projectivity, $X$ embeds over $S$ as a locally closed subscheme of
  $\PP(V)$ for some coherent sheaf $V$ on $S$, with
  $L\cong \restr{\Oo_{\PP(V)}(1)}{X}$. Let $\overline X\subset \PP(V)$ be the
  schematic closure of $X$; it is projective over $S$ and contains $X$ as an open
  subscheme. By \cref{lem:coherent-prolongation} the coherent sheaf $E$ admits a
  coherent prolongation $\overline E$ to $\overline X$. For any such prolongation,
  \cref{lem:quot-open-subfunctor} shows $\Quotf_{E/X/S}$ is an open subfunctor of
  $\Quotf_{\overline E/\overline X/S}$. The latter is representable by a
  projective $S$-scheme by \cref{thm:grothendieck-quot}, so its open subfunctor
  $\Quotf_{E/X/S}^{\Phi,L}$ is representable by a quasi-projective $S$-scheme.
\end{proof}

\section{The Hilbert scheme}

The Hilbert functor is the special case $E = \Oo_X$ of the Quot functor; its
representing scheme is the Hilbert scheme. Its existence is therefore an
immediate corollary of the main theorem.

\begin{definition}\leanok
[Hilbert scheme]
  \label{def:hilbert-scheme}
  \lean{MR2223407ConstructionHilbertQuot.hilbertScheme}
  \uses{def:hilbert-functor, def:quot-functor}
  \textit{Source: Nitsure §6 (Hilbert scheme as Quot of the structure sheaf).}
  For $X$ over $S$, a relatively very ample $L$ and $\Phi\in\QQ[\lambda]$, the
  \emph{Hilbert functor} $\Hilbf_{X/S}^{\Phi,L}$ is by definition the Quot
  functor $\Quotf_{\Oo_X/X/S}^{\Phi,L}$ of the structure sheaf: its $T$-points are
  the closed subschemes $Z\subset X_T$ flat and proper over $T$ with Hilbert
  polynomial $\Phi$ relative to $L$. When it is representable, the representing
  scheme $\hilbsch_{X/S}^{\Phi,L} = \quotsch_{\Oo_X/X/S}^{\Phi,L}$ is the
  \emph{Hilbert scheme} of $X$ over $S$.
\end{definition}

\begin{theorem}\leanok
[Existence of the Hilbert scheme]
  \label{thm:hilbert-scheme-exists}
  \lean{MR2223407ConstructionHilbertQuot.hilbertScheme_exists}
  \uses{def:hilbert-scheme, def:hilbert-functor, thm:altman-kleiman-quot, def:represents-functor}
  \textit{Source: Nitsure §6 (the $E=\Oo_X$ case of the Quot theorem).}
  Let $S$ be noetherian, $X$ strongly projective over $S$ with relatively very
  ample $L$, and $\Phi\in\QQ[\lambda]$. Then $\Hilbf_{X/S}^{\Phi,L}$ is
  representable by a scheme $\hilbsch_{X/S}^{\Phi,L}$ which is a closed subscheme
  of a suitable $\PP(F)$ over $S$; in particular the Hilbert scheme is projective
  over $S$.
\end{theorem}

\begin{proof}
  \uses{def:hilbert-scheme, thm:altman-kleiman-quot}
  By \cref{def:hilbert-scheme}, $\Hilbf_{X/S}^{\Phi,L} = \Quotf_{\Oo_X/X/S}^{\Phi,L}$.
  Apply \cref{thm:altman-kleiman-quot} with the coherent sheaf $E = \Oo_X$: it
  represents this Quot functor by a scheme embeddable as a closed subscheme of a
  projective bundle $\PP(F)$ over $S$. This representing scheme is, by definition,
  $\hilbsch_{X/S}^{\Phi,L}$.
\end{proof}

\section{Scheme of morphisms}

We first record the standard flatness facts, then the openness of the
flat-locus / isomorphism-locus, and finally the representability of the
morphisms functor as an open subscheme of a Hilbert scheme of graphs.

\begin{lemma}\leanok
[Local criterion of flatness, anchors]
  \label{lem:flatness-local-criterion}
  \lean{MR2223407ConstructionHilbertQuot.flatness_local_criterion}
  % SOURCE: Nitsure, §6, lines 2754--2771 (read from references/MR2223407-construction-hilbert-quot.tex)
  % SOURCE QUOTE: "\begin{lemma}\label{local criterion}
  % \textbf{(1) } Any finite-type flat morphism between noetherian schemes is open.
  % \textbf{(2) } Let $\pi : Y \to X$ be a finite-type morphism of
  % noetherian schemes. Then all $y\in Y$ such that $\pi$ is
  % flat at $y$ (that is, ${\mathcal O}_{Y,y}$ is a flat ${\mathcal O}_{X,\pi(y)}$-module)
  % form an open subset of $Y$.
  % \textbf{(3) } Let $S$ be a noetherian scheme, and let $f:X \to S$
  % and $g: Y\to S$ be finite type flat morphisms.
  % Let $\pi : Y \to X$ be any morphism such that $g = f\circ \pi$.
  % Let $y\in Y$, let $x=\pi(y)$, and let
  % $s = g(y) = f(x)$. If the restricted
  % morphism $\pi_s : Y_s \to X_s$ between the
  % fibers over $s$ is flat at $y\in Y_s$, then $\pi$ is flat at $y\in Y$.
  % \end{lemma}"
  \textit{Source: Nitsure §6, Lemma (local criterion); Altman--Kleiman [A-K~1] Ch.~V.}
  Let all schemes be noetherian. \textbf{(1)} A finite-type flat morphism is open.
  \textbf{(2)} For a finite-type morphism $\pi\colon Y\to X$, the set of $y\in Y$
  at which $\pi$ is flat (i.e.\ $\Oo_{Y,y}$ is flat over $\Oo_{X,\pi(y)}$) is open
  in $Y$. \textbf{(3)} Let $f\colon X\to S$ and $g\colon Y\to S$ be finite-type
  flat, and $\pi\colon Y\to X$ with $g = f\circ\pi$. For $y\in Y$, $x = \pi(y)$,
  $s = g(y)$: if the fibre morphism $\pi_s\colon Y_s\to X_s$ is flat at $y$, then
  $\pi$ is flat at $y$.
\end{lemma}

% SOURCE QUOTE PROOF: "\proof See for example Altman and Kleiman [A-K 1] Chapter V.
% The statement \textbf{(3) } is a consequence of what is known as the
% \textbf{local criterion for flatness}. \qed"
\begin{proof}
  These are standard facts; see Altman--Kleiman [A-K~1] Chapter~V. Statement
  \textbf{(3)} is the local criterion of flatness. Treated as foundational
  inputs (citations to the literature; not re-proved here).
\end{proof}

\begin{theorem}\leanok
[Flat-locus and isomorphism-locus are open]
  \label{thm:isomorphism-is-open}
  \lean{MR2223407ConstructionHilbertQuot.isomorphism_is_open}
  \uses{lem:flatness-local-criterion, thm:flattening-stratification}
  % SOURCE: Nitsure, §6, lines 2778--2794 (read from references/MR2223407-construction-hilbert-quot.tex)
  % SOURCE QUOTE: "\begin{theorem}\label{isomorphism is open}
  % Let $S$ be a noetherian scheme, and let $f:X \to S$
  % and $g: Y\to S$ be proper flat morphisms.
  % Let $\pi : Y \to X$ be any projective morphism with
  % $g = f\circ \pi$. Then $S$ has open subschemes
  % $S_2\subset S_1 \subset S$ with the following universal properties:
  % \textbf{(a) } For any locally noetherian $S$-scheme $T$, the base change
  % $\pi_T : Y_T \to X_T$ is a flat morphism if and only if the
  % structure morphism $T\to S$ factors via $S_1$. (This does not need
  % $\pi$ to be projective.)
  % \textbf{(b) } For any locally noetherian $S$-scheme $T$, the base change
  % $\pi_T : Y_T \to X_T$ is an isomorphism if and only if the
  % structure morphism $T\to S$ factors via $S_2$.
  % \end{theorem}"
  \textit{Source: Nitsure §6, Theorem (isomorphism is open).}
  Let $S$ be noetherian, $f\colon X\to S$ and $g\colon Y\to S$ proper flat, and
  $\pi\colon Y\to X$ projective with $g = f\circ\pi$. Then $S$ has open
  subschemes $S_2\subset S_1\subset S$ with the universal properties: for any
  locally noetherian $S$-scheme $T$, \textbf{(a)} the base change
  $\pi_T\colon Y_T\to X_T$ is flat iff $T\to S$ factors through $S_1$ (projectivity
  of $\pi$ not needed); \textbf{(b)} $\pi_T$ is an isomorphism iff $T\to S$
  factors through $S_2$.
\end{theorem}

% SOURCE QUOTE PROOF: "\proof \textbf{(a) } By Lemma \ref{local criterion}.(2), all
% $y\in Y$ such that $\pi$ is flat at $y$ form an
% open subset $Y'\subset Y$. Then $S_1 = S - g(Y - Y')$
% is an open subset of $S$ as $g$ is proper.
% We give $S_1$ the open subscheme structure induced from $S$.
% It follows from the local criterion of flatness
% (Lemma  \ref{local criterion}.(3))
% that $S_1$ exactly consists of all $s\in S$ such that
% the restricted morphism $\pi_s : Y_s \to X_s$ between the
% fibers over $s$ is flat. Therefore again by the
% local criterion of flatness, $S_1$ has the desired universal property.
% \textbf{(b) } Let $\pi_1 : Y_1 \to X_1$ be the pull-back of
% $\pi$ under the inclusion $S_1 \hra S$.
% Let $L$ be a relatively very ample
% line bundle for the projective morphism $\pi_1 : Y_1\to X_1$.
% Then by noetherianness there exists an integer $m\ge 1$ such that
% ${\pi_1}_*L^m$ is generated by its global sections and $R^i{\pi_1}_*L^m =0$
% for all $i \ge 1$. By flatness of $\pi_1$, it follows that
% ${\pi_1}_*L^m$ is a locally free sheaf. Let $U\subset X_1$ be the
% open subschemes such that
% ${\pi_1}_*L^m$ is of rank $1$ on $U$.
% Finally, let $S_2 = S_1 - f(X_1 - U)$, which is open as
% $f$ is proper.
% We give $S_2$ the induced open subscheme structure,
% and leave it to the reader to verify that
% it indeed has the required universal property \textbf{(b) }. \qed"
\begin{proof}
  \uses{lem:flatness-local-criterion, thm:flattening-stratification}
  \textbf{(a)} By \cref{lem:flatness-local-criterion}\,(2) the flat locus
  $Y'\subset Y$ of $\pi$ is open; since $g$ is proper, $S_1 = S\setminus g(Y\setminus Y')$
  is open, given the induced open subscheme structure. By
  \cref{lem:flatness-local-criterion}\,(3), $S_1$ consists exactly of the $s\in S$
  for which the fibre morphism $\pi_s\colon Y_s\to X_s$ is flat, and the local
  criterion gives the stated universal property (this is the flattening behaviour
  recorded in \cref{thm:flattening-stratification}).

  \textbf{(b)} Pull $\pi$ back along $S_1\hookrightarrow S$ to
  $\pi_1\colon Y_1\to X_1$ and choose a relatively very ample $L$ for $\pi_1$. By
  noetherianity there is $m\ge 1$ with $\pi_{1*}L^m$ globally generated and
  $R^i\pi_{1*}L^m = 0$ for $i\ge 1$; flatness of $\pi_1$ makes $\pi_{1*}L^m$
  locally free. Let $U\subset X_1$ be the open locus where $\pi_{1*}L^m$ has rank
  $1$, and set $S_2 = S_1\setminus f(X_1\setminus U)$, open since $f$ is proper.
  With the induced structure, $S_2$ has the universal property \textbf{(b)}.
\end{proof}

If $X,Y$ are $S$-schemes, an \emph{$S$-morphism from $X$ to $Y$ parametrised by
$T$} is a $T$-morphism $X\times_S T\to Y\times_S T$; the set of these is
$\Mor_S(X,Y)(T)$, and $T\mapsto \Mor_S(X,Y)(T)$ is a contravariant functor from
$S$-schemes to sets.

\begin{theorem}\leanok
[Scheme of morphisms]
  \label{thm:scheme-of-morphisms}
  \lean{MR2223407ConstructionHilbertQuot.schemeOfMorphisms}
  \uses{thm:hilbert-scheme-exists, thm:isomorphism-is-open, def:represents-functor, def:mor-functor}
  % SOURCE: Nitsure, §6, lines 2853--2859 (read from references/MR2223407-construction-hilbert-quot.tex)
  % SOURCE QUOTE: "\begin{theorem}
  % Let $S$ be a noetherian scheme, let $X$ be a projective scheme over $S$,
  % and let $Y$ be quasi-projective scheme
  % over $S$. Assume moreover that $X$ is flat over $S$. Then the
  % functor $\Mor_S(X,Y)$ is representable by an open
  % subscheme $\mor_S(X,Y)$ of $\hilb_{X\times_SY/S}$.
  % \end{theorem}"
  \textit{Source: Nitsure §6, Theorem (scheme of morphisms).}
  Let $S$ be noetherian, $X$ projective and flat over $S$, and $Y$
  quasi-projective over $S$. Then $\Mor_S(X,Y)$ is representable by an open
  subscheme $\Mor_S(X,Y)$ of the Hilbert scheme $\hilbsch_{X\times_S Y/S}$ of
  graphs.
\end{theorem}

% SOURCE QUOTE PROOF: "\proof We can associate
% to each morphism $f : X_T \to Y_T$ (where $T$ is a scheme over $S$)
% its graph $\Gamma_T(f) \subset (X\times_SY)_T$, which is
% closed in $(X\times_SY)_T$ by separatedness of $Y\to S$.
% We regard $\Gamma_T(f)$ as a closed subscheme of $(X\times_SY)_T$
% which is isomorphic to $X$ under the graph morphism
% $(\id_X,f): X \to  \Gamma_T(f)$ and the projection $\Gamma_T(f) \to X$,
% which are inverses to each other.
% As $X$ is proper and flat over $S$, so is $\Gamma_T(f)$, therefore
% this defines a set-map $\Gamma_T : \Mor_S(X,Y)(T) \to
% \Hilb_{X\times_SY/S}(T)$ which is functorial in $T$, so
% we obtain a morphism of functors
% $$\Gamma : \Mor_S(X,Y) \to \Hilb_{X\times_SY/S}$$
% Given any element of $\Hilb_{X\times_SY/S}(T)$, represented by a
% family $Z \subset (X\times_SY)_T$, it follows by applying
% Theorem \ref{isomorphism is open}.(b) to
% the projection $Z\to X$ that $T$ has an open subscheme
% $T'$ with the following universal property: for any
% base-change $U \to T$, the pull-back $Z_U\subset (X\times_SY)_U$
% maps isomorphically on to $X_U$ under the projection
% $p : (X\times_SY)_U\to X_U$ if and only if $U\to T$ factors via
% $T'$. Note therefore that over $T'$,
% the scheme $Z_{T'}$ will be the graph of a uniquely determined
% morphism $X_{T'} \to Y_{T'}$.
% This shows that the morphism of functors
% $\Gamma : \Mor_S(X,Y) \to \Hilb_{X\times_SY/S}$
% is a representable morphism which is an open embedding.
% Therefore a representing scheme
% $\mor_S(X,Y)$ for $\Mor_S(X,Y)$ exists as an open
% subscheme of $\hilb_{X\times_SY/S}$. \qed"
\begin{proof}
  \uses{thm:hilbert-scheme-exists, thm:isomorphism-is-open, def:represents-functor, def:mor-functor}
  To a morphism $f\colon X_T\to Y_T$ associate its graph
  $\Gamma_T(f)\subset (X\times_S Y)_T$, closed by separatedness of $Y\to S$ and
  isomorphic to $X$ via $(\id_X,f)$ and the projection. Since $X$ is proper and
  flat over $S$, so is $\Gamma_T(f)$, giving a functorial map
  $\Gamma_T\colon \Mor_S(X,Y)(T)\to \Hilbf_{X\times_S Y/S}(T)$ and hence a
  morphism of functors $\Gamma\colon \Mor_S(X,Y)\to \Hilbf_{X\times_S Y/S}$, the
  Hilbert functor being representable by \cref{thm:hilbert-scheme-exists}. Given a
  family $Z\subset (X\times_S Y)_T$, applying \cref{thm:isomorphism-is-open}\,(b)
  to the projection $Z\to X$ produces an open $T'\subset T$ over which (and only
  over which, after base change) $Z$ maps isomorphically to $X$; over $T'$, $Z$ is
  the graph of a unique morphism $X_{T'}\to Y_{T'}$. Thus $\Gamma$ is a
  representable open embedding, and $\Mor_S(X,Y)$ is represented by an open
  subscheme of $\hilbsch_{X\times_S Y/S}$.
\end{proof}

\section{Quotient by a flat projective equivalence relation}

A \emph{schematic equivalence relation} on an $S$-scheme $X$ is an $S$-scheme $R$
with a morphism $f\colon R\to X\times_S X$ such that for every $S$-scheme $T$ the
map $f(T)\colon R(T)\to X(T)\times X(T)$ is injective with image the graph of an
equivalence relation on $X(T)$. A morphism $q\colon X\to Q$ is a \emph{quotient}
for $f$ if it is a co-equaliser of the two projections
$f_1,f_2\colon R\rightrightarrows X$; the quotient is \emph{effective} if
$(f_1,f_2)\colon R\to X\times_Q X$ is an isomorphism. The construction uses the
following descent fact.

\begin{lemma}\leanok
[Descent of closed subschemes]
  \label{lem:descent-of-subschemes}
  \lean{MR2223407ConstructionHilbertQuot.descent_of_subschemes}
  \uses{thm:ffdescent}
  % SOURCE: Nitsure, §6, lines 3024--3040 (read from references/MR2223407-construction-hilbert-quot.tex)
  % SOURCE QUOTE: "\begin{lemma}\label{descent of subschemes}
  % (1) Any faithfully flat quasi-compact morphism of schemes $f : X \to Y$
  % is an effective epimorphism, that is, $f$ is a co-equaliser for
  % the projections $p_1,p_2: X\times_YX   \stackrel{\to}{\scriptstyle\to} X$.
  % (2) Let $p : D\to H$ be a faithfully flat quasi-compact morphism.
  % Let $Z \subset D$ be a closed subscheme
  % such that
  % $$p_1^{-1}Z = p_2^{-1}Z \subset D\times_HD$$
  % where $p_1,p_2 :  D\times_HD \stackrel{\to}{\scriptstyle\to} D$
  % are the projections, and $p_i^{-1}Z$ is the schematic
  % inverse image of $Z$ under $p_i$.
  % Then there exists a unique closed subscheme $Q$ of $H$
  % such that $Z = p^{-1}Q \subset D$. By base-change from $p : D\to H$,
  % it follows that the induced morphism $p|_Z : Z\to Q$
  % is faithfully flat and quasi-compact.\hfill $\square$
  % \end{lemma}"
  \textit{Source: Nitsure §6, Lemma (descent of subschemes); [SGA~1] Exp.~VIII Cor.~1.9.}
  \textbf{(1)} Any faithfully flat quasi-compact morphism $f\colon X\to Y$ is an
  effective epimorphism: $f$ is a co-equaliser of the projections
  $p_1,p_2\colon X\times_Y X\rightrightarrows X$. \textbf{(2)} Let
  $p\colon D\to H$ be faithfully flat quasi-compact and $Z\subset D$ a closed
  subscheme with $p_1^{-1}Z = p_2^{-1}Z\subset D\times_H D$. Then there is a
  unique closed subscheme $Q\subset H$ with $Z = p^{-1}Q$, and $p|_Z\colon Z\to Q$
  is faithfully flat and quasi-compact.
\end{lemma}

\begin{proof}
  \uses{thm:ffdescent}
  \leanok
  A special case of faithfully flat descent for quasi-coherent sheaves and their
  subsheaves of ideals (\cref{thm:ffdescent}): the descent datum on the ideal
  $\Ii_Z\subset \Oo_D$ given by $p_1^{-1}Z = p_2^{-1}Z$ descends uniquely to an
  ideal $\Ii_Q\subset \Oo_H$, cutting out the required $Q$; the base-change
  properties are then immediate. See [SGA~1] Exp.~VIII Cor.~1.9.
\end{proof}

\begin{theorem}\leanok
[Quotient by a flat projective equivalence relation]
  \label{thm:projective-flat-quotient}
  \lean{MR2223407ConstructionHilbertQuot.projectiveFlatQuotient}
  \uses{thm:hilbert-scheme-exists, thm:isomorphism-is-open, lem:descent-of-subschemes, thm:quasiprojective-quot}
  % SOURCE: Nitsure, §6, lines 3059--3071 (read from references/MR2223407-construction-hilbert-quot.tex)
  % SOURCE QUOTE: "\begin{theorem}\label{Projective flat quotient}
  % Let $S$ be a noetherian scheme, and let
  % $X\to S$ be a quasi-projective morphism.
  % Let $f : R \to X\times_SX$ be a schematic
  % equivalence relation on $X$ over $S$, such that
  % the projections $f_1,f_2 : R \stackrel{\to}{\scriptstyle\to} X$
  % are proper and flat.
  % Then a schematic quotient $X\to Q$ exists over $S$.
  % Moreover, $Q$ is quasi-projective over $S$,
  % the morphism $X \to Q$ is faithfully flat and projective, and
  % the induced morphism $(f_1,f_2) : R\to X\times_QX$ is an isomorphism (the
  % quotient is effective).
  % \end{theorem}"
  \textit{Source: Nitsure §6, Theorem (Projective flat quotient).}
  Let $S$ be noetherian and $X\to S$ quasi-projective. Let
  $f\colon R\to X\times_S X$ be a schematic equivalence relation whose projections
  $f_1,f_2\colon R\rightrightarrows X$ are proper and flat. Then a schematic
  quotient $X\to Q$ exists over $S$; $Q$ is quasi-projective over $S$, the
  morphism $X\to Q$ is faithfully flat and projective, and
  $(f_1,f_2)\colon R\to X\times_Q X$ is an isomorphism (the quotient is
  effective).
\end{theorem}

% SOURCE QUOTE PROOF: "\proof (Following Altman and Kleiman [A-K 2])
% The properness of $f_i$
% together with separatedness of $X\to S$ implies properness of
% $f : R \to X\times_SX$.
% Also, $f$ is functorially injective by definition
% of a schematic equivalence relation.
% It follows that $f$ is a closed embedding, which
% allows us to regard $R$ as a closed subscheme of $X\times_SX$
% (\textbf{Exercise}:
% \textit{Any proper morphism of noetherian schemes, which is injective
% at the level of functor of points, is a closed embedding}).
% This defines an element $(R)$ of $\Hilb_{X/S}(X)$, as the projection
% $p_2|_R = f_2$ is proper and flat.
% By Theorem \ref{Quasi-projective quot},
% there exists a scheme $\hilb_{X/S}$ which represents the
% functor $\Hilb_{X/S}$.
% As the parameter scheme $X$ is noetherian and as the Hilbert polynomial
% is locally constant,
% only finitely many polynomials $\Phi$ occur as Hilbert
% polynomials of fibers of $f_2: R \to X$ with respect to a
% chosen relatively very ample line bundle $L$ on $X$ over $S$.
% Let $H$ be the finite disjoint union of the corresponding
% open subschemes
% $\hilb_{X/S}^{\Phi,L}$ of $\hilb_{X/S}$.
% Then $H$ is a quasi-projective scheme over $S$ as
% each $\hilb_{X/S}^{\Phi,L}$ is so by
% Theorem \ref{Quasi-projective quot}.
% The family $(R) \in  \Hilb_{X/S}(X)$
% therefore defines a classifying morphism
% ${\varphi}: X \to H$, with the property that $\varphi^*D = R$
% where $D\subset X\times_SH$ denotes the restriction to $H$
% of the universal family over $\hilb_{X/S}$,
% and $\varphi^*D$ denotes $(\id_X\times\,{\varphi})^{-1} D$.
% Also, note that the projection $p: D\to H$ is proper and flat.
% If $X$ is non-empty then each fiber of $f_2 : R \to X$
% is also non-empty as the diagonal $\Delta_X$ is contained in $R$,
% and therefore the Hilbert polynomial of each fiber is non-zero.
% Hence $p: D\to H$ is surjective, and so $p$ is faithfully flat."
\begin{proof}
  \uses{thm:hilbert-scheme-exists, thm:isomorphism-is-open, lem:descent-of-subschemes, thm:quasiprojective-quot}
  % The full source proof (lines 3074--3278) is long; the key construction and
  % the crucial property (***) are quoted verbatim above and continued in prose.
  \emph{Following Altman--Kleiman [A-K~2].} Properness of the $f_i$ and
  separatedness of $X\to S$ make $f\colon R\to X\times_S X$ proper; being
  functorially injective, $f$ is a closed embedding, so $R$ is a closed subscheme
  of $X\times_S X$ with $p_2|_R = f_2$ proper and flat. This makes $R$ an element
  of $\Hilbf_{X/S}(X)$.

  By \cref{thm:quasiprojective-quot} (applied to $E = \Oo_X$, i.e.\
  \cref{thm:hilbert-scheme-exists}) the Hilbert functor is representable; since
  $X$ is noetherian and the Hilbert polynomial is locally constant, only finitely
  many $\Phi$ occur as Hilbert polynomials of the fibres of $f_2$ relative to a
  chosen very ample $L$. Let $H$ be the finite disjoint union of the corresponding
  $\hilbsch_{X/S}^{\Phi,L}$, a quasi-projective $S$-scheme. The family $(R)$ then
  gives a classifying morphism $\varphi\colon X\to H$ with $\varphi^*D = R$, where
  $D\subset X\times_S H$ is the universal family; the projection $p\colon D\to H$
  is proper and flat, and surjective (as $\Delta_X\subset R$ makes every fibre of
  $f_2$ non-empty), hence faithfully flat.

  One establishes the crucial property: for any $S$-scheme $T$ and
  $x,y\in X(T)$, $(x,y)\in R(T) \iff x\in\varphi(y) \iff \varphi(x)=\varphi(y)$.
  The graph $(\id_X,\varphi)\colon X\to X\times_S H$ is a closed embedding
  factoring through $D$, giving a closed subscheme
  $\Gamma_\varphi\subset D$ isomorphic to $X$. Using the property above one checks
  $p_1^{-1}\Gamma_\varphi = p_2^{-1}\Gamma_\varphi\subset D\times_H D$, so by
  \cref{lem:descent-of-subschemes}\,(2) there is a unique closed subscheme
  $Q\subset H$ with $\Gamma_\varphi = p^{-1}Q$. Define $q\colon X\to Q$ as the
  composite $X\xrightarrow{(\id_X,\varphi)}\Gamma_\varphi\xrightarrow{p}Q$, so that
  $X\xrightarrow{q}Q\hookrightarrow H$ equals $\varphi$.

  Then: \textbf{(i)} $Q$ is quasi-projective over $S$, being closed in the
  quasi-projective $H$. \textbf{(ii)} $q$ is faithfully flat and projective, by
  base change from $p\colon D\to H$ along the Cartesian square with vertices
  $X,D,Q,H$. \textbf{(iii)} By the crucial property,
  $q\circ f_1 = q\circ f_2$ and $(f_1,f_2)\colon R\to X\times_Q X$ is an
  isomorphism (the quotient is effective); under it the $f_i$ become the
  projections $p_1,p_2\colon X\times_Q X\rightrightarrows X$, and by
  \cref{lem:descent-of-subschemes}\,(1) $q$ is their co-equaliser, hence a
  co-equaliser of $f_1,f_2$. Thus $q\colon X\to Q$ is the desired quotient.

  The isomorphism-openness input \cref{thm:isomorphism-is-open} is what guarantees
  that the classifying data and the graph $\Gamma_\varphi$ are well-defined family
  members over the relevant bases.
\end{proof}

\section{Auxiliary declarations}

\begin{definition}

\leanok
[Representability by objectwise bijection]
  \label{def:represents-functor}
  \lean{MR2223407ConstructionHilbertQuot.RepresentsFunctor}
  For a presheaf $F\colon(\Sch/S)^{\mathrm{op}}\to\mathrm{Sets}$ and an object
  $Q$, the proposition that $Q$ \emph{represents} $F$: there is a family of
  bijections $\Mor_S(T,Q)\cong F(T)$, one for each test scheme $T$ (naturality in
  $T$ understood). This is the universe-polymorphic representability predicate
  used throughout the §6 corollaries; it accommodates both the
  $\mathrm{Type}\,(u+1)$-valued Quot/Hilbert functors and the $\mathrm{Type}\,u$
  functor $\Mor_S(X,Y)$ of \cref{def:mor-functor}, which the
  $\cong h_Q$ form of \cref{def:points-functor} cannot compare directly because
  of the universe gap. (A future refactor may merge the two representability
  shapes; both are honest and the §6 statements are faithful under either.)
\end{definition}

\begin{definition}

\leanok
[The functor of $S$-morphisms $\Mor_S(X,Y)$]
  \label{def:mor-functor}
  \lean{MR2223407ConstructionHilbertQuot.morFunctor}
  The contravariant functor sending a test scheme $T$ to the set of
  $T$-morphisms $X_T\to Y_T$ between the base changes
  $X_T = X\times_S T$ and $Y_T = Y\times_S T$, i.e. the functor whose
  representability is asserted by \cref{thm:scheme-of-morphisms}. The object map
  is given by the fibre products $X\times_S T$, $Y\times_S T$; the functorial
  action on morphisms is the pullback of $T$-morphisms along base change.
\end{definition}
